% Ch17 self-study -- "I do / You do" ladder for Zumdahl chapter 17.
% FIFTEEN ladders (60 questions) plus a ten-question mixed set = 70
% questions, fifteen figures.
%
% SCOPE: nothing skipped, and -- verified against the CED PDF on
% 2026-08-23 -- CED topics 9.1 through 9.7 carry NO exclusion
% statements at all. This is the only large chapter in the series with
% a completely clean scope: everything Zumdahl 17.1-17.10 teaches that
% maps to Unit 9 is assessable as written.
%
% Every number here was recomputed in Python (see the fact-check
% script), including the CaCO3 crossover temperature, which is checked
% against the real decomposition temperature of limestone.
\documentclass{apchem-worksheet}

\apcunitword{Chapter}
\apcunit{17}
\apcunittitle{Spontaneity, Entropy, and Free Energy}
\apcdoctitle{Self-Study \textbullet\ Chapter 17, I~do / You~do}
\apcsubtitle{Zumdahl \S17.1--17.10 \textbullet\ PDF pp.\ 835--869}

\begin{document}
\apctitleblock

\begin{apcnote}[How to use these notes --- read this first]
Each skill is a \textbf{ladder}: a fully worked example in the
solid-framed box, then \textbf{four} YOUR TURN questions in the dashed
box --- same skill, new numbers, no help. Work all four before checking
the gray \emph{check:} line; a tracker tick needs all four right first
time. A ten-question \textbf{mixed set} follows Ladder 15.

\textbf{Nothing in this chapter is skipped, and nothing is excluded.}
CED topics 9.1 to 9.7 carry no exclusion statements at all --- unusual,
and worth knowing. Where other chapters have material you can safely
skim, this one does not.

\textbf{The question this chapter answers.} Chapter 12 asked
\emph{how fast}. Chapter 13 asked \emph{how far}. This chapter asks the
one underneath both: \textbf{why does anything happen at all?} Why does
ice melt in a warm room and never freeze in one? Why do gases mix and
never unmix?

\textbf{The vocabulary warning that will cost you marks if ignored.}
The CED says \textbf{``thermodynamically favoured''}, not
``spontaneous''. Zumdahl says spontaneous. They mean the same thing ---
$\Delta G < 0$ --- but ``spontaneous'' misleads, because it sounds like
``happens immediately'' and a thermodynamically favoured reaction can
take geological time (Ladder 11). Use the CED's phrase in answers.
\end{apcnote}

% =====================================================================
\section*{\sffamily Ladder 1 \textbullet\ What entropy actually is}
\zsec{17.1}

Entropy, $S$, measures how many ways the energy and matter of a system
can be arranged. More available arrangements means higher entropy.

\begin{center}
\begin{tikzpicture}[font=\sffamily\footnotesize]
  % confined gas
  \node[font=\sffamily\small\bfseries] at (1.9,3.3) {confined};
  \fill[apcgray!22] (0.6,0.5) rectangle (3.2,2.7);
  \draw[very thick] (0.6,2.7) rectangle (3.2,0.5);
  \draw[very thick,dashed] (1.9,0.5) -- (1.9,2.7);
  \foreach \p in {(0.95,0.9),(1.35,1.5),(0.85,2.1),(1.55,2.35),
                  (1.15,1.15),(1.6,0.75)} {
    \fill \p circle (0.075);}
  \node[align=center] at (1.9,-0.15)
    {all particles on one side\\\textbf{few} arrangements --- low $S$};

  \draw[-{Stealth[length=3mm]},very thick] (3.7,1.6) -- (5.0,1.6);
  \node[above,font=\sffamily\footnotesize] at (4.35,1.75) {remove};
  \node[below,font=\sffamily\footnotesize] at (4.35,1.45) {the barrier};

  % spread gas
  \begin{scope}[xshift=5.5cm]
    \node[font=\sffamily\small\bfseries] at (1.9,3.3) {spread out};
    \fill[apcgray!22] (0.6,0.5) rectangle (3.2,2.7);
    \draw[very thick] (0.6,2.7) rectangle (3.2,0.5);
    \foreach \p in {(0.95,1.1),(1.7,1.9),(2.4,0.85),(2.9,1.7),
                    (1.3,2.35),(2.15,2.3)} {
      \fill \p circle (0.075);}
    \node[align=center] at (1.9,-0.15)
      {particles anywhere\\\textbf{many} arrangements --- high $S$};
  \end{scope}
\end{tikzpicture}
\end{center}

\begin{apctrap}
\textbf{``Entropy is disorder'' is a slogan, not a definition, and it
misleads.} A crystal growing from solution looks more ordered, yet the
total entropy still rises --- because the heat released disperses into
the surroundings. Say \textbf{``the number of ways energy and matter
can be arranged''} or \textbf{``dispersal of energy''}. Answers built on
``messiness'' lose marks the moment a question involves the
surroundings.
\end{apctrap}

\begin{workedexample}[I do: why the gas never goes back]
\textbf{Remove the barrier above and the gas fills the container. Wait
a year and it never returns to one side. Why not?}

\textbf{Nothing forbids it.} No law of motion prevents every molecule
from happening to be on the left at the same instant. The molecules do
not know where the barrier used to be.

\textbf{It is overwhelmingly improbable.} For one molecule the chance
of being on the left is $1/2$. For $n$ molecules all being on the left
at once, it is $(1/2)^{n}$. For a mole --- $6 \times 10^{23}$ molecules
--- that probability is zero for every practical purpose.

\textbf{So entropy is a counting statement.} The spread-out
arrangement is not \emph{preferred} by any force; there are simply
vastly more ways to be spread out than to be bunched. Systems drift
toward the arrangement with the most ways of happening, and with $10^{23}$
particles that drift is irreversible.

\textbf{Three consequences worth carrying.} Gases expanding into a
vacuum, two gases mixing, and heat flowing from hot to cold are all the
same phenomenon: energy and matter spreading over more available
arrangements. None of the three requires a force pushing them; each is
just the more numerous outcome winning.
\end{workedexample}

\begin{yourturn}[YOUR TURN 1 --- four questions]
\begin{enumerate}[leftmargin=*,itemsep=0.5em]
  \item Define entropy without using the word ``disorder''.
        \workspace[1.6cm]
  \item Two gases in connected flasks mix and never unmix. Explain
        using arrangements. \workspace[1.8cm]
  \item Does anything physically \emph{forbid} the gas returning to one
        side? Explain. \workspace[1.8cm]
  \item Which has higher entropy: \SI{1}{\mole} of \ce{Ar} in
        \SI{1}{\litre} or in \SI{10}{\litre}? Why?
        \workspace[1.6cm]
\end{enumerate}
\selfcheck{(a) the number of ways energy and matter can be arranged
\quad (b) far more mixed arrangements exist \quad (c) no --- it is
merely overwhelmingly improbable \quad (d) \SI{10}{\litre} --- more
positions available}
\keyonly{\begin{apcanswer}(a) The \textbf{number of ways the energy and
matter of a system can be arranged} --- equivalently, how far energy is
dispersed. (b) There are \textbf{vastly more arrangements} in which the
two gases are intermingled than in which they are separated, so the
mixed state is overwhelmingly the more probable one; nothing pulls them
together again. (c) \textbf{No.} No physical law forbids it --- it is
simply so improbable ($(1/2)^{n}$ for $n$ molecules) that with a mole of
gas it never happens. (d) \textbf{\SI{10}{\litre}}. A larger volume
offers each atom far more available positions, so there are more
arrangements and $S$ is higher.\end{apcanswer}}
\end{yourturn}

% =====================================================================
\section*{\sffamily Ladder 2 \textbullet\ Predicting the sign of
$\Delta S$}
\zsec{17.1}

You can call the sign of $\Delta S$ without a single number, and the
exam asks you to constantly.

\begin{center}
\begin{tikzpicture}[font=\sffamily\footnotesize]
  \node[font=\sffamily\small\bfseries] at (5.6,3.5)
    {check these in order --- the first that applies decides};
  \foreach \y/\n/\t in {%
      2.85/{1}/{\textbf{moles of gas} change? more gas $=$ higher $S$
        --- this dominates everything},
      2.3/{2}/{\textbf{phase} change? solid $<$ liquid $\ll$ gas},
      1.75/{3}/{\textbf{dissolving} a solid? usually higher $S$
        (Ladder 14)},
      1.2/{4}/{\textbf{temperature} rising? always higher $S$},
      0.65/{5}/{more \textbf{particles} on one side? usually higher
        $S$}} {
    \node[anchor=west,font=\sffamily\bfseries] at (0.5,\y) {\n.};
    \node[anchor=west] at (1.05,\y) {\t};}
  \draw[apcgray] (0.4,3.2) -- (11.0,3.2);
  \draw[apcgray] (0.4,0.3) -- (11.0,0.3);
\end{tikzpicture}
\end{center}

\begin{workedexample}[I do: four calls, no calculations]
\textbf{Give the sign of $\Delta S$ for each.}

\textbf{(a) $\ce{2NO2(g) -> N2O4(g)}$.} Two moles of gas become one.
Fewer gas particles means fewer arrangements: $\Delta S$ is
\textbf{negative}. Counting gas moles is the first thing to check and
usually the last thing you need.

\textbf{(b) $\ce{H2O(l) -> H2O(g)}$.} A liquid becomes a gas ---
molecules go from touching to flying free. Strongly \textbf{positive},
and this is the single biggest entropy change in ordinary chemistry.

\textbf{(c) $\ce{CaCO3(s) -> CaO(s) + CO2(g)}$.} No gas on the left,
one mole of gas on the right. \textbf{Positive} --- and note that the
two solids barely matter. Creating gas from no gas is decisive.

\textbf{(d) $\ce{2H2(g) + O2(g) -> 2H2O(l)}$.} Three moles of gas
become zero. Strongly \textbf{negative}.

\textbf{The hierarchy is what makes this fast.} If the number of gas
moles changes, that answers the question --- stop there. Only when gas
moles are equal on both sides do you go looking at phases, particle
counts or dissolving. Students who weigh all five factors equally get
slow and get it wrong.

\textbf{And when gas moles are equal, expect a small $\Delta S$.} For
$\ce{H2(g) + I2(g) -> 2HI(g)}$, two moles of gas become two: whatever
$\Delta S$ turns out to be, it will be near zero, and you should be
suspicious of any calculation that says otherwise.
\end{workedexample}

\begin{yourturn}[YOUR TURN 2 --- four questions]
Give the sign of $\Delta S$ and one-line reasoning.
\begin{enumerate}[leftmargin=*,itemsep=0.5em]
  \item $\ce{2SO2(g) + O2(g) -> 2SO3(g)}$ \workspace[1.6cm]
  \item $\ce{H2O(s) -> H2O(l)}$ \workspace[1.6cm]
  \item $\ce{NH4NO3(s) -> NH4+(aq) + NO3-(aq)}$
        \workspace[1.6cm]
  \item $\ce{N2(g) + 3H2(g) -> 2NH3(g)}$ \workspace[1.6cm]
\end{enumerate}
\selfcheck{(a) negative --- 3 gas mol to 2 \quad (b) positive --- solid
to liquid \quad (c) positive --- solid to dissolved ions \quad (d)
negative --- 4 gas mol to 2}
\keyonly{\begin{apcanswer}(a) \textbf{Negative}: three moles of gas
become two. (b) \textbf{Positive}: melting frees molecules from fixed
lattice positions. (c) \textbf{Positive}: an ordered solid becomes two
kinds of freely moving hydrated ion --- and note this is the entropy
that drives cold packs (Ladder 14). (d) \textbf{Negative}: four moles of
gas become two.\end{apcanswer}}
\end{yourturn}

% =====================================================================
\section*{\sffamily Ladder 3 \textbullet\ The third law and absolute
entropies}
\zsec{17.2}

Entropy has something enthalpy does not: a true zero. That makes $S$
values \textbf{absolute}, not relative.

\begin{center}
\begin{tikzpicture}[font=\sffamily\footnotesize]
  \draw[-{Stealth[length=3mm]},very thick] (0.8,0.6) -- (0.8,3.4)
    node[above,font=\sffamily\footnotesize] {$S$};
  \draw[very thick] (0.65,0.9) -- (0.95,0.9);
  \node[anchor=east,font=\sffamily\bfseries] at (0.55,0.9) {0};
  \node[anchor=west] at (1.15,0.9)
    {a \textbf{perfect crystal at \SI{0}{\kelvin}} --- one arrangement
     only};
  \node[anchor=west] at (1.15,1.7)
    {every real substance sits \emph{above} zero};
  \node[anchor=west] at (1.15,2.5)
    {so tabulated $S^{\circ}$ values are \textbf{positive}, even for
     elements};
  % Separator sits BELOW the body, before the closing note. At y=3.05
  % it was above every line and read as a stray top border.
  \draw[apcgray] (0.6,0.45) -- (10.9,0.45);
  \node[anchor=west,font=\sffamily\bfseries] at (0.8,-0.05)
    {contrast: $\Delta H_{f}^{\circ}$ of an element is \emph{defined}
     as 0 --- $S^{\circ}$ of an element is \emph{measured} and is not};
\end{tikzpicture}
\end{center}

\begin{apctrap}
\textbf{$S^{\circ}$ for an element is not zero.} $\Delta H_{f}^{\circ}$
of \ce{O2(g)} is zero by definition; $S^{\circ}$ of \ce{O2(g)} is
\SI{205.1}{\joule\per\mole\per\kelvin}, a measured quantity. Setting
element entropies to zero in a $\Delta S^{\circ}$ calculation is the
standard wrong answer in Ladder 4, and it comes from carrying a
Chapter 6 habit into Chapter 17.
\end{apctrap}

\begin{workedexample}[I do: reading an $S^{\circ}$ table]
\textbf{Three patterns you can predict before looking anything up.}

\textbf{Phase dominates.} For water: $S^{\circ}(\text{s}) = 41$,
$S^{\circ}(\text{l}) = 69.9$, $S^{\circ}(\text{g}) =
\SI{188.8}{\joule\per\mole\per\kelvin}$. The gas is nearly three times
the liquid. Whenever you compare two substances, check the phases
first.

\textbf{Bigger and more complex means higher.} $S^{\circ}(\ce{CH4}) =
186$ but $S^{\circ}(\ce{C2H6}) =
\SI{229.5}{\joule\per\mole\per\kelvin}$. More atoms means more ways to
vibrate and rotate, so more arrangements of the same energy.

\textbf{Harder solids have lower entropy.} Diamond sits at
\SI{2.4}{\joule\per\mole\per\kelvin} --- almost nothing --- because its
rigid covalent network leaves atoms very few ways to move. Graphite,
with its sliding sheets, is \SI{5.7}{}. The Chapter 10 structure
argument reappears as a number.

\textbf{Why this matters for the third law.} A perfect crystal at
\SI{0}{\kelvin} has exactly one arrangement, so $S = 0$ exactly. Every
other state has more, so every real $S^{\circ}$ is positive. That is
what makes entropy absolute --- and it is why you never see a negative
$S^{\circ}$ in a table, though $\Delta S^{\circ}$ for a reaction is
negative all the time.
\end{workedexample}

\begin{yourturn}[YOUR TURN 3 --- four questions]
\begin{enumerate}[leftmargin=*,itemsep=0.5em]
  \item State the third law of thermodynamics. \workspace[1.6cm]
  \item Why is $S^{\circ}(\ce{O2}) \neq 0$ when
        $\Delta H_{f}^{\circ}(\ce{O2}) = 0$? \workspace[1.8cm]
  \item Which is larger, $S^{\circ}(\ce{H2O(l)})$ or
        $S^{\circ}(\ce{H2O(g)})$? \workspace[1.4cm]
  \item Which is larger, $S^{\circ}(\ce{CH4})$ or
        $S^{\circ}(\ce{C3H8})$? Why? \workspace[1.6cm]
\end{enumerate}
\selfcheck{(a) a perfect crystal at \SI{0}{\kelvin} has $S = 0$ \quad
(b) $S$ is absolute; $\Delta H_{f}$ is a defined reference \quad (c)
the gas \quad (d) \ce{C3H8} --- more atoms, more arrangements}
\keyonly{\begin{apcanswer}(a) The entropy of a \textbf{perfect crystal
at \SI{0}{\kelvin} is zero} --- there is exactly one way to arrange it.
(b) Because entropy has a \textbf{true zero} and is measured
\emph{absolutely} from it, whereas $\Delta H_{f}^{\circ}$ is a
\emph{defined reference point} --- elements are assigned zero by
convention, not by measurement. (c) \textbf{The gas}, by a large margin
(188.8 against \SI{69.9}{\joule\per\mole\per\kelvin}). (d)
\textbf{\ce{C3H8}}: more atoms give more ways to vibrate and rotate, so
more arrangements at a given energy.\end{apcanswer}}
\end{yourturn}

% =====================================================================
\section*{\sffamily Ladder 4 \textbullet\ Calculating
$\Delta S^{\circ}_{\text{rxn}}$}
\zsec{17.6}

Same shape as Hess's law from Chapter 6 --- products minus reactants,
each scaled by its coefficient.

\[ \Delta S^{\circ}_{\text{rxn}}
   = \sum n\,S^{\circ}(\text{products})
   - \sum n\,S^{\circ}(\text{reactants}) \]

\begin{workedexample}[I do: a calculation whose sign you already knew]
\textbf{For $\ce{2H2(g) + O2(g) -> 2H2O(l)}$, find
$\Delta S^{\circ}$. Use $S^{\circ}$: \ce{H2(g)} $= 130.7$,
\ce{O2(g)} $= 205.1$, \ce{H2O(l)} $=
\SI{69.9}{\joule\per\mole\per\kelvin}$.}

\textbf{Predict the sign first.} Three moles of gas become zero moles
of gas. $\Delta S^{\circ}$ must be strongly negative --- so if the
arithmetic comes out positive, the arithmetic is wrong.

\textbf{Products:}
\[ 2 \times 69.9 = \SI{139.8}{\joule\per\kelvin} \]

\textbf{Reactants:}
\[ 2 \times 130.7 + 205.1 = 261.4 + 205.1
   = \SI{466.5}{\joule\per\kelvin} \]

\textbf{Subtract:}
\[ \Delta S^{\circ} = 139.8 - 466.5
   = \mathbf{\SI{-326.7}{\joule\per\kelvin}} \]

Negative, as predicted, and large --- consistent with destroying three
moles of gas.

\textbf{Two habits.} Predict the sign before you compute; it costs five
seconds and catches every subtraction reversal. And watch the
\textbf{units}: $S^{\circ}$ is tabulated in
\si{\joule\per\mole\per\kelvin} while $\Delta H$ is in
\si{\kilo\joule\per\mole}. Ladder 7 mixes the two in one equation, and
that unit mismatch is the most expensive error in this chapter.

\textbf{Do not set element entropies to zero.} \ce{H2} and \ce{O2} are
elements, but their $S^{\circ}$ values are 130.7 and 205.1, not 0.
Dropping them here would give $\Delta S^{\circ} =
\SI{+139.8}{\joule\per\kelvin}$ --- positive, for a reaction that
destroys three moles of gas. The sign check catches it instantly.
\end{workedexample}

\begin{yourturn}[YOUR TURN 4 --- four questions]
Use $S^{\circ}$ (\si{\joule\per\mole\per\kelvin}): \ce{N2(g)} $= 191.6$,
\ce{H2(g)} $= 130.7$, \ce{NH3(g)} $= 192.5$, \ce{CaCO3(s)} $= 92.9$,
\ce{CaO(s)} $= 39.8$, \ce{CO2(g)} $= 213.7$.
\begin{enumerate}[leftmargin=*,itemsep=0.5em]
  \item Predict the sign of $\Delta S^{\circ}$ for
        $\ce{N2 + 3H2 -> 2NH3}$. \workspace[1.4cm]
  \item Compute it. \workspace[2.0cm]
  \item Compute $\Delta S^{\circ}$ for
        $\ce{CaCO3(s) -> CaO(s) + CO2(g)}$. \workspace[2.0cm]
  \item Why may you not treat $S^{\circ}(\ce{N2})$ as zero?
        \workspace[1.6cm]
\end{enumerate}
\selfcheck{(a) negative \quad (b) \SI{-198.7}{\joule\per\kelvin} \quad
(c) \SI{+160.6}{\joule\per\kelvin} \quad (d) $S^{\circ}$ is absolute,
not a defined reference}
\keyonly{\begin{apcanswer}(a) \textbf{Negative} --- four moles of gas
become two. (b) Products $2 \times 192.5 = 385.0$; reactants
$191.6 + 3 \times 130.7 = 191.6 + 392.1 = 583.7$; so
$\Delta S^{\circ} = 385.0 - 583.7 =
\mathbf{\SI{-198.7}{\joule\per\kelvin}}$ --- negative, as predicted.
(c) Products $39.8 + 213.7 = 253.5$; reactant $92.9$; so
$\Delta S^{\circ} = 253.5 - 92.9 =
\mathbf{\SI{+160.6}{\joule\per\kelvin}}$ --- positive, because a gas is
created from none. (d) Because $S^{\circ}$ values are \textbf{absolute},
measured from the third-law zero. Only $\Delta H_{f}^{\circ}$ uses
elements as a defined zero.\end{apcanswer}}
\end{yourturn}

% =====================================================================
\section*{\sffamily Ladder 5 \textbullet\ The second law:
$\Delta S_{\text{universe}}$}
\zsec{17.2}

The actual criterion for a favoured process is not about the system at
all.

\[ \Delta S_{\text{univ}}
   = \Delta S_{\text{sys}} + \Delta S_{\text{surr}} > 0
   \quad\text{for a thermodynamically favoured process} \]

\begin{center}
\begin{tikzpicture}[font=\sffamily\footnotesize]
  \fill[apcgray!12] (0.5,0.4) rectangle (7.4,3.1);
  \draw[very thick] (0.5,0.4) rectangle (7.4,3.1);
  \node[anchor=north west,font=\sffamily\bfseries] at (0.7,2.95)
    {the universe};
  \fill[apcgray!35] (1.4,0.9) rectangle (4.2,2.3);
  \draw[very thick] (1.4,0.9) rectangle (4.2,2.3);
  \node[font=\sffamily\bfseries] at (2.8,1.9) {the system};
  \node at (2.8,1.35) {$\Delta S_{\text{sys}}$};
  \node[font=\sffamily\bfseries] at (5.9,1.9) {the surroundings};
  \node at (5.9,1.35) {$\Delta S_{\text{surr}}$};
  \draw[{Stealth[length=2mm]}-{Stealth[length=2mm]},thick]
    (4.35,1.6) -- (4.95,1.6);
  \node[anchor=west,align=left] at (8.0,1.9)
    {only the \textbf{sum} must};
  \node[anchor=west,align=left] at (8.0,1.5)
    {be positive --- either};
  \node[anchor=west,align=left] at (8.0,1.1)
    {piece may be negative};
\end{tikzpicture}
\end{center}

\begin{workedexample}[I do: how order can increase]
\textbf{Water freezing at \SI{-10}{\celsius} is thermodynamically
favoured. Yet liquid becomes solid --- $\Delta S_{\text{sys}}$ is
clearly negative. How can that be?}

\textbf{Because the system is not the universe.} Freezing releases heat
--- the enthalpy of fusion, outward --- and that heat disperses into the
surrounding air. Dispersing energy into the surroundings \emph{raises}
$\Delta S_{\text{surr}}$.

\textbf{Below \SI{0}{\celsius} the surroundings win.} The heat released
raises $\Delta S_{\text{surr}}$ by more than $\Delta S_{\text{sys}}$
falls, so $\Delta S_{\text{univ}} > 0$ and freezing is favoured.

\textbf{Above \SI{0}{\celsius} it does not.} Ladder 6 shows why:
$\Delta S_{\text{surr}} = -\Delta H/T$, so the same released heat buys
\emph{less} entropy in warmer surroundings. At high enough $T$ the
surroundings' gain no longer covers the system's loss, and melting ---
the reverse --- becomes the favoured direction instead.

\textbf{So \SI{0}{\celsius} is not a magic temperature.} It is simply
where the two terms exactly cancel, $\Delta S_{\text{univ}} = 0$, and
neither direction is favoured. That is the definition of a phase
equilibrium, and it is the same idea as $\Delta G = 0$ in Ladder 7.

\textbf{The exam version of this idea:} whenever a question shows an
apparently order-increasing process happening on its own --- freezing,
crystallising, a protein folding --- the answer involves heat released
to the surroundings.
\end{workedexample}

\begin{yourturn}[YOUR TURN 5 --- four questions]
\begin{enumerate}[leftmargin=*,itemsep=0.5em]
  \item State the second law in terms of $\Delta S_{\text{univ}}$.
        \workspace[1.6cm]
  \item Can a favoured process have $\Delta S_{\text{sys}} < 0$?
        Explain. \workspace[1.8cm]
  \item Water freezes at \SI{-10}{\celsius}. Give the signs of
        $\Delta S_{\text{sys}}$ and $\Delta S_{\text{surr}}$.
        \workspace[1.8cm]
  \item What is $\Delta S_{\text{univ}}$ for a system exactly at
        equilibrium? \workspace[1.4cm]
\end{enumerate}
\selfcheck{(a) $\Delta S_{\text{univ}} > 0$ for a favoured process
\quad (b) yes, if $\Delta S_{\text{surr}}$ is more positive \quad (c)
sys negative, surr positive \quad (d) zero}
\keyonly{\begin{apcanswer}(a) For any thermodynamically favoured
process, $\Delta S_{\text{univ}} > 0$; the entropy of the universe
increases. (b) \textbf{Yes.} Only the sum must be positive, so a
system's entropy may fall provided the surroundings gain more --- which
happens whenever the process releases enough heat. (c)
$\Delta S_{\text{sys}}$ is \textbf{negative} (liquid to solid);
$\Delta S_{\text{surr}}$ is \textbf{positive} (heat released disperses
into the surroundings) and larger in magnitude. (d)
\textbf{Zero} --- neither direction is favoured, which is what
equilibrium means.\end{apcanswer}}
\end{yourturn}

% =====================================================================
\section*{\sffamily Ladder 6 \textbullet\ $\Delta S_{\text{surr}}$ from
$\Delta H$}
\zsec{17.2}

The surroundings only ever see the heat the system gives or takes, so
their entropy change follows one small equation.

\[ \Delta S_{\text{surr}} = \frac{-\Delta H_{\text{sys}}}{T} \]

\begin{center}
\begin{tikzpicture}[font=\sffamily\footnotesize]
  \node[font=\sffamily\small\bfseries] at (5.5,2.9)
    {read the equation twice --- the sign, then the temperature};
  \node[anchor=west] at (0.5,2.25)
    {\textbf{exothermic} ($\Delta H < 0$): heat flows \emph{out}, so
     $\Delta S_{\text{surr}}$ is \textbf{positive}};
  \node[anchor=west] at (0.5,1.7)
    {\textbf{endothermic} ($\Delta H > 0$): heat flows \emph{in}, so
     $\Delta S_{\text{surr}}$ is \textbf{negative}};
  \draw[apcgray] (0.4,1.35) -- (11.0,1.35);
  \node[anchor=west] at (0.5,0.85)
    {and the $T$ in the denominator: the \emph{same} heat matters
     \textbf{more} to cold surroundings};
  \node[anchor=west,font=\sffamily\itshape] at (0.5,0.3)
    {a coin means more to a beggar than to a millionaire --- same
     energy, different context};
\end{tikzpicture}
\end{center}

\begin{workedexample}[I do: the temperature in the denominator]
\textbf{The Haber reaction has $\Delta H =
\SI{-92.2}{\kilo\joule\per\mole}$. Find $\Delta S_{\text{surr}}$ at
\SI{298}{\kelvin}.}

\textbf{Convert to joules first} --- $\Delta S$ is always in joules per
kelvin, never kilojoules:
\[ \Delta S_{\text{surr}}
   = \frac{-(\num{-92200}\;\si{\joule\per\mole})}{\SI{298}{\kelvin}}
   = \mathbf{\SI{+309}{\joule\per\mole\per\kelvin}} \]

Positive, because an exothermic reaction dumps heat into the
surroundings and spreads it there.

\textbf{Now the same reaction at \SI{700}{\kelvin}:}
\[ \Delta S_{\text{surr}} = \frac{92200}{700}
   = \SI{+132}{\joule\per\mole\per\kelvin} \]

\textbf{Less than half, from identical heat.} That is the whole content
of the $T$ in the denominator: hot surroundings are already full of
dispersed energy, so a given quantity of heat adds proportionally fewer
new arrangements.

\textbf{And this is why exothermic reactions stop being favoured when
heated.} $\Delta S_{\text{sys}}$ does not care much about temperature,
but $\Delta S_{\text{surr}}$ shrinks as $T$ rises. Eventually it can no
longer carry a negative $\Delta S_{\text{sys}}$, and the reaction
reverses --- exactly the Haber-process compromise from Chapter 13,
Ladder 12, now derived rather than asserted.

\textbf{This equation is also where $\Delta G$ comes from.} Multiply
$\Delta S_{\text{univ}} = \Delta S_{\text{sys}} - \Delta H/T$ through by
$-T$ and you get $-T\Delta S_{\text{univ}} = \Delta H - T\Delta S$ ---
the right-hand side is $\Delta G$. Ladder 7 uses it; it is not a new
idea, just a rearrangement.
\end{workedexample}

\begin{yourturn}[YOUR TURN 6 --- four questions]
\begin{enumerate}[leftmargin=*,itemsep=0.5em]
  \item A reaction has $\Delta H = \SI{-50.0}{\kilo\joule\per\mole}$.
        Find $\Delta S_{\text{surr}}$ at \SI{250}{\kelvin}.
        \workspace[1.8cm]
  \item Same reaction at \SI{500}{\kelvin}. \workspace[1.8cm]
  \item Give the sign of $\Delta S_{\text{surr}}$ for an endothermic
        process. \workspace[1.4cm]
  \item Why does the same released heat produce a larger
        $\Delta S_{\text{surr}}$ at low temperature?
        \workspace[1.8cm]
\end{enumerate}
\selfcheck{(a) \SI{+200}{\joule\per\mole\per\kelvin} \quad (b)
\SI{+100}{\joule\per\mole\per\kelvin} \quad (c) negative \quad (d) cold
surroundings hold less dispersed energy already}
\keyonly{\begin{apcanswer}(a) $-(\num{-50000})/250 =
\mathbf{\SI{+200}{\joule\per\mole\per\kelvin}}$. (b) $50000/500 =
\mathbf{\SI{+100}{\joule\per\mole\per\kelvin}}$ --- exactly half, since
the temperature doubled. (c) \textbf{Negative}: an endothermic process
draws heat \emph{out} of the surroundings, reducing their entropy. (d)
Because cold surroundings already contain relatively little dispersed
energy, so a given quantity of heat represents a proportionally
\textbf{larger} increase in the number of available arrangements. The
$T$ in the denominator says exactly this.\end{apcanswer}}
\end{yourturn}

% =====================================================================
\section*{\sffamily Ladder 7 \textbullet\ Gibbs free energy}
\zsec{17.3--17.4}

$\Delta G$ bundles the system and the surroundings into a single
number you can compute from system properties alone.

\[ \Delta G = \Delta H - T\,\Delta S \]

\begin{center}
\begin{tikzpicture}[font=\sffamily\footnotesize]
  \draw[-{Stealth[length=3mm]},very thick] (0.6,2.0) -- (10.9,2.0);
  \node[font=\sffamily\bfseries] at (5.7,2.5) {$\Delta G = 0$};
  \fill (5.7,2.0) circle (0.11);
  \draw[densely dashed,apcgray] (5.7,1.72) -- (5.7,1.95);
  \node[font=\sffamily\bfseries] at (2.6,1.55) {$\Delta G < 0$};
  \node[align=center] at (2.6,0.9)
    {thermodynamically\\\textbf{favoured}};
  \node[align=center] at (5.7,0.9) {at\\equilibrium};
  \node[font=\sffamily\bfseries] at (8.8,1.55) {$\Delta G > 0$};
  \node[align=center] at (8.8,0.9)
    {\textbf{not} favoured\\(the reverse is)};
  \draw[apcgray] (0.5,0.35) -- (11.0,0.35);
  \node[anchor=west] at (0.6,-0.1)
    {$\Delta H$ in \si{\kilo\joule}, $\Delta S$ in
     \si{\joule\per\kelvin} --- \textbf{convert one} before
     subtracting};
\end{tikzpicture}
\end{center}

\begin{apctrap}
\textbf{The unit mismatch is the single most expensive error in this
chapter.} $\Delta H$ is tabulated in \textbf{kilojoules};
$\Delta S$ in \textbf{joules} per kelvin. In $\Delta H - T\Delta S$ you
must convert one to match the other. Forget, and $T\Delta S$ comes out a
thousand times too large --- which usually flips the sign of $\Delta G$
and therefore the entire conclusion.
\end{apctrap}

\begin{workedexample}[I do: is limestone decomposition favoured at room
temperature?]
\textbf{For $\ce{CaCO3(s) -> CaO(s) + CO2(g)}$, $\Delta H^{\circ} =
\SI{+178.3}{\kilo\joule}$ and $\Delta S^{\circ} =
\SI{+160.6}{\joule\per\kelvin}$. Find $\Delta G^{\circ}$ at
\SI{298}{\kelvin}.}

\textbf{Convert $\Delta S$ to kilojoules} so both terms match:
\[ \Delta S^{\circ} = \SI{160.6}{\joule\per\kelvin}
   = \SI{0.1606}{\kilo\joule\per\kelvin} \]

\textbf{Substitute:}
\[ \Delta G^{\circ} = 178.3 - (298)(0.1606)
   = 178.3 - 47.9 = \mathbf{\SI{+130.4}{\kilo\joule}} \]

\textbf{Positive, so not favoured at \SI{298}{\kelvin}} --- limestone
does not spontaneously crumble into quicklime on a shelf, which
matches every cathedral still standing.

\textbf{But look at the structure of the answer.} $\Delta H$ is
positive (unfavourable) and $\Delta S$ is positive (favourable). They
oppose each other, and at \SI{298}{\kelvin} the enthalpy term is nearly
four times larger, so enthalpy wins.

\textbf{Raise the temperature and that changes.} $T\Delta S$ grows in
proportion to $T$ while $\Delta H$ barely moves --- so at some
temperature the entropy term overtakes and $\Delta G$ turns negative.
Ladder 9 finds exactly where. This is why lime kilns are heated: not to
make the reaction faster, but to make it \emph{favoured at all}.
\end{workedexample}

\begin{yourturn}[YOUR TURN 7 --- four questions]
\begin{enumerate}[leftmargin=*,itemsep=0.5em]
  \item $\Delta H = \SI{-100.0}{\kilo\joule}$, $\Delta S =
        \SI{+50.0}{\joule\per\kelvin}$, $T = \SI{298}{\kelvin}$. Find
        $\Delta G$. \workspace[1.8cm]
  \item $\Delta H = \SI{+40.0}{\kilo\joule}$, $\Delta S =
        \SI{+100.0}{\joule\per\kelvin}$, $T = \SI{298}{\kelvin}$. Find
        $\Delta G$ and say whether it is favoured. \workspace[2.0cm]
  \item What does $\Delta G = 0$ mean? \workspace[1.4cm]
  \item Before computing, why must you convert units?
        \workspace[1.6cm]
\end{enumerate}
\selfcheck{(a) \SI{-114.9}{\kilo\joule} \quad (b)
\SI{+10.2}{\kilo\joule}, not favoured \quad (c) equilibrium \quad (d)
$\Delta H$ is in kJ, $\Delta S$ in J}
\keyonly{\begin{apcanswer}(a) $\Delta G = -100.0 - (298)(0.0500) =
-100.0 - 14.9 = \mathbf{\SI{-114.9}{\kilo\joule}}$ --- favoured, with
both terms helping. (b) $\Delta G = 40.0 - (298)(0.1000) = 40.0 - 29.8
= \mathbf{\SI{+10.2}{\kilo\joule}}$: \textbf{not favoured} at this
temperature, though it would be at higher $T$ (Ladder 9). (c) The system
is at \textbf{equilibrium} --- neither direction is favoured. (d)
Because $\Delta H$ is tabulated in \textbf{kilojoules} and $\Delta S$ in
\textbf{joules} per kelvin; subtracting without converting makes the
$T\Delta S$ term a thousand times too large and usually flips the sign
of the answer.\end{apcanswer}}
\end{yourturn}

% =====================================================================
\section*{\sffamily Ladder 8 \textbullet\ The four sign combinations}
\zsec{17.4}

Two signs, four cases, and only two of them depend on temperature.

\begin{center}
\begin{tikzpicture}[font=\sffamily\footnotesize]
  \node[font=\sffamily\small\bfseries] at (5.9,3.7)
    {when is $\Delta G = \Delta H - T\Delta S$ negative?};
  \node[font=\sffamily\bfseries] at (4.3,3.05) {$\Delta S > 0$};
  \node[font=\sffamily\bfseries] at (8.3,3.05) {$\Delta S < 0$};
  \node[anchor=east,font=\sffamily\bfseries] at (2.5,2.25)
    {$\Delta H < 0$};
  \node[anchor=east,font=\sffamily\bfseries] at (2.5,1.15)
    {$\Delta H > 0$};
  \draw[apcgray] (2.7,2.7) -- (10.2,2.7);
  \draw[apcgray] (2.7,1.65) -- (10.2,1.65);
  \draw[apcgray] (6.3,3.3) -- (6.3,0.55);
  \node[align=center] at (4.3,2.25)
    {\textbf{always favoured}\\at every $T$};
  \node[align=center] at (8.3,2.25)
    {favoured at\\\textbf{low} $T$};
  \node[align=center] at (4.3,1.15)
    {favoured at\\\textbf{high} $T$};
  \node[align=center] at (8.3,1.15)
    {\textbf{never} favoured\\at any $T$};
  \node[font=\sffamily\footnotesize] at (5.9,0.1)
    {the two diagonal boxes are temperature-dependent --- those are the
     ones Ladder 9 is for};
\end{tikzpicture}
\end{center}

\begin{workedexample}[I do: reading the table without memorising it]
\textbf{You never need to memorise these four boxes. Read them off
$\Delta G = \Delta H - T\Delta S$ directly.}

\textbf{$\Delta H < 0$, $\Delta S > 0$.} The first term is negative and
$-T\Delta S$ is also negative. Both terms push $\Delta G$ negative, so
it is negative at every temperature. \emph{Example:} combustion.

\textbf{$\Delta H > 0$, $\Delta S < 0$.} Both terms are positive, so
$\Delta G$ is positive at every temperature. Nothing you do to $T$ can
help. \emph{Example:} the reverse of combustion --- which is why fires
do not un-burn.

\textbf{$\Delta H < 0$, $\Delta S < 0$.} Enthalpy helps, entropy hurts,
and the entropy term grows with $T$. So it is favoured while $T$ is
\textbf{small} enough. \emph{Example:} freezing, or the Haber process.

\textbf{$\Delta H > 0$, $\Delta S > 0$.} Enthalpy hurts, entropy helps
and grows with $T$. Favoured once $T$ is \textbf{large} enough.
\emph{Example:} limestone decomposition (Ladder 7), or ice melting.

\textbf{The one-line rule that regenerates the whole table:}
\emph{$T$ multiplies only the entropy term, so raising the temperature
always strengthens whatever entropy is arguing for.} If $\Delta S$ is
positive, heat helps; if negative, heat hurts. Everything else follows.
\end{workedexample}

\begin{yourturn}[YOUR TURN 8 --- four questions]
\begin{enumerate}[leftmargin=*,itemsep=0.5em]
  \item $\Delta H < 0$, $\Delta S > 0$: at what temperatures is it
        favoured? \workspace[1.4cm]
  \item $\Delta H > 0$, $\Delta S > 0$: at what temperatures?
        \workspace[1.4cm]
  \item A reaction is favoured only below \SI{400}{\kelvin}. Give the
        signs of $\Delta H$ and $\Delta S$. \workspace[1.6cm]
  \item Why does raising $T$ always favour whatever $\Delta S$ is
        arguing for? \workspace[1.8cm]
\end{enumerate}
\selfcheck{(a) all \quad (b) high $T$ only \quad (c) both negative
\quad (d) $T$ multiplies only the entropy term}
\keyonly{\begin{apcanswer}(a) \textbf{All temperatures} --- both terms
push $\Delta G$ negative. (b) \textbf{High temperatures only}: the
positive $\Delta H$ must be overcome by $T\Delta S$, which requires a
large $T$. (c) Both \textbf{negative}: enthalpy favours it, entropy
opposes it, so it succeeds only while $T\Delta S$ is small --- that is,
at low temperature. (d) Because in $\Delta G = \Delta H - T\Delta S$,
temperature multiplies \textbf{only} the entropy term. Raising $T$
magnifies $\Delta S$'s contribution without changing $\Delta H$'s, so
whatever entropy wants, heat gives it more of.\end{apcanswer}}
\end{yourturn}

% =====================================================================
\section*{\sffamily Ladder 9 \textbullet\ The crossover temperature}
\zsec{17.4}

For the two temperature-dependent cases, there is one temperature where
$\Delta G$ passes through zero. Set $\Delta G = 0$ and solve.

\[ 0 = \Delta H - T\Delta S
   \qquad\Longrightarrow\qquad
   T = \frac{\Delta H}{\Delta S} \]

\begin{center}
\begin{tikzpicture}
\begin{axis}[
  width=10.6cm, height=5.8cm,
  xlabel={temperature (\si{\kelvin})}, ylabel={$\Delta G$
    (\si{\kilo\joule})},
  xmin=0, xmax=1600, ymin=-90, ymax=200,
  xtick={0,400,800,1110,1600},
  xticklabels={0,400,800,1110,1600},
  ytick={-50,0,50,100,150},
  grid=major, grid style={apcgray!25},
  label style={font=\footnotesize}, tick label style={font=\footnotesize},
]
\addplot[seriesA, very thick, domain=0:1600, samples=2]
  {178.3 - 0.1606*x};
\addplot[black, dashed, domain=0:1600, samples=2] {0};
\addplot[only marks, mark=*, mark size=1.8pt, black]
  coordinates {(1110,0)};
\node[font=\footnotesize,anchor=west] at (axis cs:120,150)
  {$\Delta G = 178.3 - 0.1606\,T$};
% Both annotations moved clear of the descending line. The line reaches
% y=-40 at x=1359 (inside the old caption's span) and y=45 at x=830
% (inside the old "crossover" label), so each was crossed by it.
\node[font=\footnotesize,anchor=east] at (axis cs:1040,-58)
  {favoured only beyond \SI{1110}{\kelvin}};
\node[font=\footnotesize,anchor=west] at (axis cs:1150,28)
  {crossover};
\end{axis}
\end{tikzpicture}
\end{center}

\begin{workedexample}[I do: at what temperature will limestone
decompose?]
\textbf{Using $\Delta H^{\circ} = \SI{+178.3}{\kilo\joule}$ and
$\Delta S^{\circ} = \SI{+160.6}{\joule\per\kelvin}$ from Ladder 7,
find the crossover temperature.}

\textbf{Set $\Delta G = 0$ and solve, converting units as always:}
\[ T = \frac{\Delta H}{\Delta S}
     = \frac{\SI{178300}{\joule}}{\SI{160.6}{\joule\per\kelvin}}
     = \mathbf{\SI{1110}{\kelvin}} \]

That is about \SI{837}{\celsius}.

\textbf{Check it against reality.} Industrial lime kilns run at roughly
\SI{900}{\celsius} --- just above our crossover, as they must be. The
calculation predicted a real engineering constraint from two tabulated
numbers.

\textbf{Now say which side is which.} $\Delta S$ is positive, so the
entropy term grows with $T$; the reaction is therefore favoured
\emph{above} \SI{1110}{\kelvin} and not below. If $\Delta S$ had been
negative, the inequality would flip and the reaction would be favoured
\emph{below} the crossover.

\textbf{Two cautions.} First, $T = \Delta H / \Delta S$ only makes sense
when $\Delta H$ and $\Delta S$ have the \emph{same} sign --- otherwise
there is no crossover, because the reaction is favoured either always or
never (Ladder 8), and the formula returns a meaningless negative
temperature. Second, the answer is in \textbf{kelvin}, always; convert
to celsius only if the question asks.

\textbf{And a caveat worth stating on an FRQ:} this assumes
$\Delta H^{\circ}$ and $\Delta S^{\circ}$ are temperature-independent.
They are not exactly, but over ordinary ranges the approximation is
good, and it is what the exam expects.
\end{workedexample}

\begin{yourturn}[YOUR TURN 9 --- four questions]
\begin{enumerate}[leftmargin=*,itemsep=0.5em]
  \item $\Delta H = \SI{+30.0}{\kilo\joule}$, $\Delta S =
        \SI{+100.0}{\joule\per\kelvin}$. Find the crossover
        temperature. \workspace[1.8cm]
  \item Is that reaction favoured above or below it? Why?
        \workspace[1.6cm]
  \item $\Delta H = \SI{-40.0}{\kilo\joule}$, $\Delta S =
        \SI{-80.0}{\joule\per\kelvin}$. Find the crossover and say
        which side is favoured. \workspace[2.0cm]
  \item Why does $T = \Delta H/\Delta S$ give nothing useful when the
        two signs differ? \workspace[1.8cm]
\end{enumerate}
\selfcheck{(a) \SI{300}{\kelvin} \quad (b) above --- $\Delta S$ is
positive \quad (c) \SI{500}{\kelvin}, favoured below \quad (d) there is
no crossover; it is favoured always or never}
\keyonly{\begin{apcanswer}(a) $T = 30000/100.0 =
\mathbf{\SI{300}{\kelvin}}$. (b) \textbf{Above.} $\Delta S$ is positive,
so the $T\Delta S$ term grows with temperature and eventually overcomes
the positive $\Delta H$. (c) $T = \num{-40000}/(-80.0) =
\mathbf{\SI{500}{\kelvin}}$, and since $\Delta S$ is \textbf{negative}
the reaction is favoured \textbf{below} \SI{500}{\kelvin}. (d) Because
with opposite signs both terms of $\Delta G = \Delta H - T\Delta S$
point the same way at every temperature --- the reaction is favoured
either always or never, and there is no crossing point. The formula then
returns a negative temperature, which is physically
meaningless.\end{apcanswer}}
\end{yourturn}

% =====================================================================
\section*{\sffamily Ladder 10 \textbullet\ $\Delta G^{\circ}$ from
formation values}
\zsec{17.4}

A second route to $\Delta G^{\circ}$, identical in shape to Hess's law
and to Ladder 4.

\[ \Delta G^{\circ}_{\text{rxn}}
   = \sum n\,\Delta G_{f}^{\circ}(\text{products})
   - \sum n\,\Delta G_{f}^{\circ}(\text{reactants}) \]

\begin{workedexample}[I do: two routes, one answer]
\textbf{You now have two ways to reach $\Delta G^{\circ}$, and knowing
which to use is half the skill.}

\textbf{Route 1 --- $\Delta G^{\circ} = \Delta H^{\circ} -
T\Delta S^{\circ}$.} Use this when the question gives you $\Delta H$
and $\Delta S$, or when it asks about a temperature \emph{other than}
\SI{298}{\kelvin}. It is the only route that can handle a temperature
change.

\textbf{Route 2 --- formation values.} Use this when the question hands
you a $\Delta G_{f}^{\circ}$ table. It works only at
\SI{298}{\kelvin}, because that is the temperature the table was
compiled at.

\textbf{Worked: $\ce{Fe2O3(s) + 3CO(g) -> 2Fe(s) + 3CO2(g)}$}, with
$\Delta G_{f}^{\circ}$ (\si{\kilo\joule\per\mole}):
\ce{Fe2O3} $= -742.2$, \ce{CO} $= -137.2$, \ce{CO2} $= -394.4$,
\ce{Fe} $= 0$.

\textbf{Products:}
\[ 2(0) + 3(-394.4) = \SI{-1183.2}{\kilo\joule} \]

\textbf{Reactants:}
\[ -742.2 + 3(-137.2) = -742.2 - 411.6 = \SI{-1153.8}{\kilo\joule} \]

\textbf{Subtract:}
\[ \Delta G^{\circ} = -1183.2 - (-1153.8)
   = \mathbf{\SI{-29.4}{\kilo\joule}} \]

Negative, so iron smelting in a blast furnace is thermodynamically
favoured --- as several thousand years of ironworking suggest.

\textbf{Elements are zero here.} $\Delta G_{f}^{\circ}(\ce{Fe}) = 0$
because iron in its standard state is the reference --- exactly like
$\Delta H_{f}^{\circ}$ in Chapter 6. This is the \emph{opposite} of
Ladder 4's rule for $S^{\circ}$, and keeping the two straight is what
this ladder tests.
\end{workedexample}

\begin{yourturn}[YOUR TURN 10 --- four questions]
$\Delta G_{f}^{\circ}$ (\si{\kilo\joule\per\mole}): \ce{NH3(g)} $=
-16.4$, \ce{NO(g)} $= +86.6$, \ce{H2O(g)} $= -228.6$, elements $= 0$.
\begin{enumerate}[leftmargin=*,itemsep=0.5em]
  \item Find $\Delta G^{\circ}$ for $\ce{N2(g) + 3H2(g) -> 2NH3(g)}$.
        \workspace[1.8cm]
  \item Find $\Delta G^{\circ}$ for $\ce{N2(g) + O2(g) -> 2NO(g)}$.
        \workspace[1.8cm]
  \item Which of the two is thermodynamically favoured at
        \SI{298}{\kelvin}? \workspace[1.4cm]
  \item Why is $\Delta G_{f}^{\circ}$ of an element zero when
        $S^{\circ}$ of an element is not? \workspace[1.8cm]
\end{enumerate}
\selfcheck{(a) \SI{-32.8}{\kilo\joule} \quad (b)
\SI{+173.2}{\kilo\joule} \quad (c) the first \quad (d)
$\Delta G_{f}$ is a defined reference; $S$ is absolute}
\keyonly{\begin{apcanswer}(a) $2(-16.4) - [0 + 3(0)] =
\mathbf{\SI{-32.8}{\kilo\joule}}$. (b) $2(+86.6) - [0 + 0] =
\mathbf{\SI{+173.2}{\kilo\joule}}$. (c) \textbf{The first}, ammonia
synthesis --- $\Delta G^{\circ}$ is negative. The second is strongly
positive, which is why atmospheric \ce{N2} and \ce{O2} coexist without
reacting (Chapter 13, Ladder 5, where $K \approx 10^{-30}$ said the same
thing). (d) Because $\Delta G_{f}^{\circ}$ is defined \emph{relative} to
elements in their standard states --- a chosen reference --- while
$S^{\circ}$ is measured \emph{absolutely} from the third-law zero. Two
different kinds of quantity that happen to appear in similar-looking
tables.\end{apcanswer}}
\end{yourturn}

% =====================================================================
\section*{\sffamily Ladder 11 \textbullet\ Thermodynamic versus kinetic
control}
\zsec{17.7}

$\Delta G$ says whether a reaction \emph{can} go. It says nothing at all
about whether it \emph{will} in your lifetime.

\begin{center}
\begin{tikzpicture}[font=\sffamily\footnotesize]
  \draw[-{Stealth[length=2.5mm]},thick] (0.6,0.4) -- (0.6,3.9)
    node[above,font=\footnotesize] {$G$};
  \draw[-{Stealth[length=2.5mm]},thick] (0.6,0.4) -- (9.6,0.4)
    node[right,font=\footnotesize] {progress};
  \draw[very thick] (1.1,2.2) -- (2.0,2.2)
    .. controls (3.0,2.2) and (3.2,3.5) .. (4.2,3.5)
    .. controls (5.2,3.5) and (5.4,1.7) .. (6.4,1.7) -- (8.6,1.7);
  \draw[dotted,thick] (1.1,2.2) -- (8.9,2.2);
  \draw[dotted,thick] (1.1,1.7) -- (8.9,1.7);
  \draw[{Stealth[length=2mm]}-{Stealth[length=2mm]},thick]
    (8.75,1.7) -- (8.75,2.2);
  \node[anchor=west,font=\footnotesize] at (8.85,1.95)
    {$\Delta G < 0$};
  \draw[{Stealth[length=2mm]}-{Stealth[length=2mm]},thick]
    (2.5,2.2) -- (2.5,3.5);
  \node[anchor=east,font=\footnotesize] at (2.4,2.85)
    {huge $E_{\mathrm{a}}$};
  \node[anchor=south,font=\footnotesize] at (1.5,2.25) {diamond};
  \node[anchor=north,font=\footnotesize] at (7.4,1.65) {graphite};
  \node[font=\sffamily\footnotesize] at (5.0,-0.15)
    {favoured, and yet it never happens --- kinetics, not
     thermodynamics, is in charge};
\end{tikzpicture}
\end{center}

\begin{workedexample}[I do: diamonds are not forever, technically]
\textbf{$\ce{C(diamond) -> C(graphite)}$ has $\Delta G^{\circ} =
\SI{-2.9}{\kilo\joule\per\mole}$. Every diamond in existence is
thermodynamically unstable. Why do jewellers sleep at night?}

\textbf{Because $\Delta G$ is not a rate.} A negative $\Delta G$ says
the products sit lower in free energy, so the conversion is favoured. It
says nothing about the barrier between them.

\textbf{And that barrier is enormous.} Converting diamond to graphite
means breaking a rigid three-dimensional network of covalent bonds and
rebuilding it as sheets. The activation energy is so large that at room
temperature the rate is immeasurably slow --- far slower than the age of
the universe.

\textbf{So the reaction is under \emph{kinetic control}.} What you
observe is decided by the barrier, not by the free-energy difference.
The diamond is not stable; it is \textbf{metastable} --- trapped behind
a hill it cannot climb.

\textbf{Two more cases of the same thing.} A hydrogen--oxygen mixture
has $\Delta G^{\circ} \approx \SI{-457}{\kilo\joule\per\mole}$ and sits
unchanged for years without a spark (Chapter 13, Ladder 5, said this in
$K$ language). Petrol in air is likewise wildly favoured and entirely
safe until ignited.

\textbf{The exam sentence.} ``The reaction is thermodynamically
favoured ($\Delta G < 0$) but kinetically controlled: the activation
energy is large, so the rate is negligible at this temperature.'' Both
halves, or you have only answered half the question.
\end{workedexample}

\begin{yourturn}[YOUR TURN 11 --- four questions]
\begin{enumerate}[leftmargin=*,itemsep=0.5em]
  \item A reaction has $\Delta G^{\circ} = \SI{-200}{\kilo\joule}$ but
        no observable change occurs. Explain. \workspace[1.8cm]
  \item Does a catalyst change $\Delta G^{\circ}$? Explain.
        \workspace[1.8cm]
  \item What does ``metastable'' mean? \workspace[1.6cm]
  \item Which quantity would you change to make the reaction in (a)
        actually proceed? \workspace[1.6cm]
\end{enumerate}
\selfcheck{(a) large $E_{\mathrm{a}}$ --- kinetic control \quad (b) no
--- it changes only the rate \quad (c) favoured but trapped behind a
barrier \quad (d) $E_{\mathrm{a}}$, via a catalyst or heat}
\keyonly{\begin{apcanswer}(a) The reaction is \textbf{thermodynamically
favoured but kinetically controlled}: the activation energy is so large
that the rate is negligible at this temperature. $\Delta G$ describes
extent, not speed. (b) \textbf{No.} A catalyst lowers $E_{\mathrm{a}}$
and so raises the rate, but $\Delta G^{\circ}$ depends only on the free
energies of reactants and products, which a catalyst does not touch ---
the same argument as Chapter 12, Ladder 12. (c) A state that is
\textbf{thermodynamically unfavoured relative to some other state but
kinetically trapped} behind a large activation barrier --- favoured to
change, but unable to. (d) The \textbf{activation energy}: add a
catalyst, or raise the temperature so more collisions clear the
existing barrier. Neither changes $\Delta G^{\circ}$.\end{apcanswer}}
\end{yourturn}

% =====================================================================
\section*{\sffamily Ladder 12 \textbullet\ $\Delta G^{\circ}$ and $K$}
\zsec{17.9}

Chapter 13's equilibrium constant and this chapter's free energy are
the same information in two currencies.

\[ \Delta G^{\circ} = -RT \ln K
   \qquad
   R = 8.314\;\mathrm{J\,mol^{-1}\,K^{-1}} \]

\begin{center}
\begin{tikzpicture}[font=\sffamily\footnotesize]
  \node[font=\sffamily\small\bfseries] at (5.6,2.9)
    {the two languages, translated};
  % NB: the loop variables are \gv/\kv/\mng, NOT \g/\k/\meaning.
  % \meaning is a TeX primitive and \k is LaTeX's ogonek accent;
  % using either as a \foreach variable is a fatal error.
  \foreach \y/\gv/\kv/\mng in {%
      2.25/{$\Delta G^{\circ} < 0$}/{$K > 1$}/{products favoured},
      1.7/{$\Delta G^{\circ} = 0$}/{$K = 1$}/{neither side favoured},
      1.15/{$\Delta G^{\circ} > 0$}/{$K < 1$}/{reactants favoured}} {
    \node[anchor=west] at (1.0,\y) {\gv};
    \node[anchor=west] at (4.2,\y) {\kv};
    \node[anchor=west] at (6.6,\y) {\mng};}
  \draw[apcgray] (0.5,2.6) -- (11.0,2.6);
  \draw[apcgray] (0.5,0.8) -- (11.0,0.8);
  \node[anchor=west] at (1.0,0.3)
    {note $\ln 1 = 0$, so $K = 1$ is exactly where
     $\Delta G^{\circ}$ changes sign};
\end{tikzpicture}
\end{center}

\begin{workedexample}[I do: both directions of the translation]
\textbf{(a) A reaction has $K = \num{1.0e6}$ at \SI{298}{\kelvin}. Find
$\Delta G^{\circ}$.}
\[ \Delta G^{\circ} = -RT \ln K
   = -(8.314)(298)\ln(\num{1.0e6}) \]
\[ = -(2478)(13.82) = \SI{-34200}{\joule}
   = \mathbf{\SI{-34.2}{\kilo\joule}} \]

Large $K$, negative $\Delta G^{\circ}$ --- the two statements agree, as
they must.

\textbf{(b) A reaction has $\Delta G^{\circ} =
\SI{+20.0}{\kilo\joule}$ at \SI{298}{\kelvin}. Find $K$.}

Rearrange:
\[ \ln K = \frac{-\Delta G^{\circ}}{RT}
         = \frac{-20000}{(8.314)(298)}
         = \frac{-20000}{2478} = -8.07 \]
\[ K = e^{-8.07} = \mathbf{\num{3.1e-4}} \]

Positive $\Delta G^{\circ}$, $K$ well below 1 --- reactants favoured,
consistent again.

\textbf{The sign check that catches everything.} Before computing,
decide which side of 1 the answer should fall on. Negative
$\Delta G^{\circ}$ gives $K > 1$; positive gives $K < 1$. A $K$ of
\num{3.1e-4} for a positive $\Delta G^{\circ}$ passes; a $K$ of 3200
would not, and would mean a dropped minus sign.

\textbf{Units, again.} $R$ is in \textbf{joules}, so $\Delta G^{\circ}$
must be in joules too. Feeding kilojoules into this equation is the same
mistake as Ladder 7's, in a new place.

\textbf{Why the relationship exists at all.} $\Delta G^{\circ}$ is the
free-energy difference between pure reactants and pure products, and $K$
measures where the mixture settles between them. They are two
descriptions of one thing --- which is why Chapter 13, Ladder 5's
``$K \gg 1$ means products favoured'' and this chapter's
``$\Delta G^{\circ} < 0$ means favoured'' were always the same
sentence.
\end{workedexample}

\begin{yourturn}[YOUR TURN 12 --- four questions]
Use $T = \SI{298}{\kelvin}$, $RT = \SI{2478}{\joule\per\mole}$.
\begin{enumerate}[leftmargin=*,itemsep=0.5em]
  \item $K = 1.0$. Find $\Delta G^{\circ}$. \workspace[1.6cm]
  \item $\Delta G^{\circ} = \SI{-10.0}{\kilo\joule}$. Is $K$ greater or
        less than 1? \workspace[1.6cm]
  \item $K = \num{1.0e-3}$. Find $\Delta G^{\circ}$.
        \workspace[2.0cm]
  \item A student computes $K = 500$ from $\Delta G^{\circ} =
        \SI{+15}{\kilo\joule}$. What went wrong?
        \workspace[1.8cm]
\end{enumerate}
\selfcheck{(a) 0 \quad (b) greater than 1 \quad (c)
\SI{+17.1}{\kilo\joule} \quad (d) sign error --- positive
$\Delta G^{\circ}$ needs $K < 1$}
\keyonly{\begin{apcanswer}(a) $\ln 1 = 0$, so $\Delta G^{\circ} =
\mathbf{0}$. (b) \textbf{Greater than 1} --- a negative
$\Delta G^{\circ}$ always corresponds to $K > 1$. (c)
$\Delta G^{\circ} = -(2478)\ln(\num{1.0e-3}) = -(2478)(-6.908) =
\SI{17120}{\joule} = \mathbf{\SI{+17.1}{\kilo\joule}}$. (d) A
\textbf{sign error}. A positive $\Delta G^{\circ}$ must give
$K < 1$; getting 500 means the minus sign in
$\ln K = -\Delta G^{\circ}/RT$ was dropped. The correct answer is
$K = e^{-15000/2478} = \num{2.3e-3}$.\end{apcanswer}}
\end{yourturn}

% =====================================================================
\section*{\sffamily Ladder 13 \textbullet\ $\Delta G$ versus
$\Delta G^{\circ}$}
\zsec{17.9}

The degree symbol matters. $\Delta G^{\circ}$ describes standard
conditions; $\Delta G$ describes the mixture you actually have.

\[ \Delta G = \Delta G^{\circ} + RT \ln Q \]

\begin{center}
\begin{tikzpicture}[font=\sffamily\footnotesize]
  \node[font=\sffamily\small\bfseries] at (5.6,2.85)
    {two symbols, two jobs};
  \node[anchor=west] at (0.6,2.2)
    {$\Delta G^{\circ}$ \quad a \textbf{fixed} number for the
     reaction at standard conditions --- tied to $K$};
  \node[anchor=west] at (0.6,1.65)
    {$\Delta G$ \quad\; \textbf{changes} as the reaction proceeds ---
     tied to $Q$};
  \draw[apcgray] (0.5,1.3) -- (11.0,1.3);
  \node[anchor=west] at (0.6,0.8)
    {at equilibrium $Q = K$ and $\Delta G = 0$ --- substitute those and
     you recover $\Delta G^{\circ} = -RT\ln K$};
  \node[anchor=west,font=\sffamily\itshape] at (0.6,0.25)
    {so Ladder 12's equation is just this one, evaluated at
     equilibrium};
\end{tikzpicture}
\end{center}

\begin{apctrap}
\textbf{$\Delta G^{\circ} = 0$ and $\Delta G = 0$ mean completely
different things.} $\Delta G = 0$ means \emph{this particular mixture}
is at equilibrium --- extremely common. $\Delta G^{\circ} = 0$ means
$K = 1$ exactly --- rare, and a much stronger claim. A reaction with
$\Delta G^{\circ} = \SI{+50}{\kilo\joule}$ still reaches equilibrium;
it just does so with very little product.
\end{apctrap}

\begin{workedexample}[I do: why a positive $\Delta G^{\circ}$ still
produces something]
\textbf{A reaction has $\Delta G^{\circ} = \SI{+20.0}{\kilo\joule}$, so
$K = \num{3.1e-4}$ (Ladder 12). Start with pure reactants and no
product. What happens?}

\textbf{At the very start, $Q = 0$.} And $\ln Q \to -\infty$, so
\[ \Delta G = \Delta G^{\circ} + RT\ln Q \]
is enormously negative. \textbf{The reaction runs forward} --- even
though $\Delta G^{\circ}$ is positive.

\textbf{As product accumulates, $Q$ rises} and $RT \ln Q$ climbs from
$-\infty$ toward zero and beyond. $\Delta G$ becomes less and less
negative.

\textbf{When $Q$ reaches $K$, $\Delta G$ hits zero} and the reaction
stops changing. Because $K$ is small, that happens with very little
product formed --- but it is not zero product.

\textbf{So the standard confusion resolves like this.} A positive
$\Delta G^{\circ}$ does \emph{not} mean nothing happens. It means the
equilibrium position lies toward reactants. Some product always forms,
because starting from $Q = 0$ the forward direction is always initially
favoured.

\textbf{And the reverse case.} A reaction with a large negative
$\Delta G^{\circ}$ started with mostly \emph{product} will run
\emph{backwards} --- because $Q > K$ makes $\Delta G$ positive. Which
direction a mixture moves is a $Q$-versus-$K$ question (Chapter 13,
Ladder 6), and this equation is that comparison written in energy
units.
\end{workedexample}

\begin{yourturn}[YOUR TURN 13 --- four questions]
\begin{enumerate}[leftmargin=*,itemsep=0.5em]
  \item State the difference between $\Delta G$ and
        $\Delta G^{\circ}$. \workspace[1.8cm]
  \item What is $\Delta G$ at equilibrium? \workspace[1.4cm]
  \item A reaction has $\Delta G^{\circ} > 0$ and starts with pure
        reactants. Does any product form? Explain.
        \workspace[1.8cm]
  \item If $Q > K$, what is the sign of $\Delta G$, and which way does
        the reaction run? \workspace[1.8cm]
\end{enumerate}
\selfcheck{(a) standard conditions vs the actual mixture \quad (b) zero
\quad (c) yes --- at $Q = 0$, $\Delta G$ is negative \quad (d)
positive; it runs in reverse}
\keyonly{\begin{apcanswer}(a) $\Delta G^{\circ}$ applies at
\textbf{standard conditions} and is a fixed number tied to $K$;
$\Delta G$ applies to \textbf{the mixture you actually have} and changes
as the reaction proceeds, tied to $Q$. (b) \textbf{Zero}. (c)
\textbf{Yes.} With pure reactants $Q = 0$, so $RT\ln Q \to -\infty$ and
$\Delta G$ is strongly negative --- the reaction runs forward until
$Q = K$. A positive $\Delta G^{\circ}$ only means the equilibrium lies
toward reactants, not that nothing happens. (d) $\Delta G$ is
\textbf{positive}, so the forward reaction is not favoured and the
system runs \textbf{in reverse} until $Q$ falls back to
$K$.\end{apcanswer}}
\end{yourturn}

% =====================================================================
\section*{\sffamily Ladder 14 \textbullet\ Free energy of dissolution}
\zsec{17.5}

Chapter 11's question --- why does anything dissolve? --- finally gets
its quantitative answer.

\begin{workedexample}[I do: the salt that dissolves while cooling the
beaker]
\textbf{\ce{NaCl} dissolving has $\Delta H =
\SI{+3.9}{\kilo\joule\per\mole}$ and $\Delta S =
\SI{+43}{\joule\per\mole\per\kelvin}$. Is it favoured at
\SI{298}{\kelvin}?}

\[ \Delta G = 3.9 - (298)(0.043)
   = 3.9 - 12.8 = \mathbf{\SI{-8.9}{\kilo\joule\per\mole}} \]

Negative, so yes --- and note \emph{how} it is negative. The enthalpy
term is unfavourable; the entropy term carries the whole reaction.

\textbf{This is the answer Chapter 11 owed you.} Ladder 5 there showed
that dissolving \ce{NaCl} is slightly endothermic --- the beaker cools
--- and asked why it happens anyway. The answer is that ions going from
a rigid lattice to freely moving hydrated particles is a large entropy
increase, and $T\Delta S$ beats a small positive $\Delta H$.

\textbf{The three-step cycle, now in free-energy terms.} Breaking the
lattice and opening the solvent both cost enthalpy; solvation returns
it. Whatever small residue is left over, the entropy term usually
overwhelms --- which is why most salts dissolve at all.

\textbf{When it fails.} For a salt with a large lattice energy and a
strongly positive $\Delta H_{\text{soln}}$, the entropy term cannot
cover it and $\Delta G > 0$ --- the salt is insoluble. That is
Chapter 16's $K_{sp} \ll 1$, and via $\Delta G^{\circ} = -RT\ln K$
it is literally the same statement: a positive $\Delta G^{\circ}$ for
dissolution \emph{is} a small $K_{sp}$.

\textbf{And the temperature prediction.} Since $\Delta S$ is usually
positive for dissolving, $T\Delta S$ grows with heat, so most solids
become more soluble when warmed --- Chapter 10, Ladder 9's solubility
curves, predicted rather than tabulated.
\end{workedexample}

\begin{yourturn}[YOUR TURN 14 --- four questions]
\begin{enumerate}[leftmargin=*,itemsep=0.5em]
  \item A salt dissolves with $\Delta H = \SI{+25.0}{\kilo\joule\per\mole}$
        and $\Delta S = \SI{+100.0}{\joule\per\mole\per\kelvin}$. Find
        $\Delta G$ at \SI{298}{\kelvin}. \workspace[2.0cm]
  \item Is it favoured? What is driving it? \workspace[1.6cm]
  \item Why is $\Delta S$ almost always positive for dissolving a
        solid? \workspace[1.8cm]
  \item A salt has $\Delta G^{\circ} > 0$ for dissolution. What does
        that say about $K_{sp}$? \workspace[1.8cm]
\end{enumerate}
\selfcheck{(a) \SI{-4.8}{\kilo\joule\per\mole} \quad (b) yes ---
entropy \quad (c) lattice to free hydrated ions \quad (d)
$K_{sp} < 1$ --- sparingly soluble}
\keyonly{\begin{apcanswer}(a) $\Delta G = 25.0 - (298)(0.1000) =
25.0 - 29.8 = \mathbf{\SI{-4.8}{\kilo\joule\per\mole}}$. (b)
\textbf{Yes, favoured}, and \textbf{entropy} is driving it --- the
enthalpy term is strongly unfavourable and is overcome by $T\Delta S$.
(c) Because ions or molecules go from \textbf{fixed positions in a rigid
lattice to moving freely through the solvent}, an enormous increase in
the number of available arrangements. (d) That $K_{sp} < 1$, and the
more positive $\Delta G^{\circ}$ is the smaller $K_{sp}$ becomes ---
by $\Delta G^{\circ} = -RT\ln K$, a positive $\Delta G^{\circ}$ for
dissolution \emph{is} a sparingly soluble salt.\end{apcanswer}}
\end{yourturn}

% =====================================================================
\section*{\sffamily Ladder 15 \textbullet\ Coupled reactions}
\zsec{17.10}

An unfavourable reaction can be driven by attaching it to a favourable
one --- provided they share a common species.

\begin{center}
\begin{tikzpicture}[font=\sffamily\footnotesize]
  \node[font=\sffamily\small\bfseries] at (5.6,3.1)
    {add the equations, add the free energies};
  \node[anchor=west] at (0.6,2.45)
    {unfavourable: \quad $\Delta G^{\circ}_{1} = +13.8$};
  \node[anchor=west] at (0.6,1.9)
    {favourable: \qquad\; $\Delta G^{\circ}_{2} = -30.5$};
  \draw[apcgray] (0.5,1.55) -- (7.6,1.55);
  \node[anchor=west,font=\sffamily\bfseries] at (0.6,1.1)
    {coupled: \qquad\quad $\Delta G^{\circ} = -16.7$
     \; \si{\kilo\joule\per\mole}};
  \node[anchor=west,align=left,font=\sffamily\footnotesize]
    at (8.2,1.9)
    {$\Delta G^{\circ}$ values \textbf{add},\\
     exactly as $\Delta H$ does\\
     in Hess's law};
\end{tikzpicture}
\end{center}

\begin{workedexample}[I do: how your cells cheat thermodynamics]
\textbf{Phosphorylating glucose is uphill:}
\[ \ce{glucose + P_i -> glucose-6-phosphate}
   \qquad \Delta G^{\circ} = \SI{+13.8}{\kilo\joule\per\mole} \]
On its own it essentially does not happen.

\textbf{ATP hydrolysis is downhill:}
\[ \ce{ATP + H2O -> ADP + P_i}
   \qquad \Delta G^{\circ} = \SI{-30.5}{\kilo\joule\per\mole} \]

\textbf{Couple them --- they share \ce{P_i}} --- and add:
\[ \ce{glucose + ATP -> glucose-6-phosphate + ADP} \]
\[ \Delta G^{\circ} = +13.8 - 30.5
   = \mathbf{\SI{-16.7}{\kilo\joule\per\mole}} \]

Negative, so the coupled reaction is favoured, and this is the first
step of glycolysis --- happening in your cells right now.

\textbf{No law was broken.} The uphill reaction did not become
favourable; it was \emph{paid for} by a larger downhill one. The
combined process still has $\Delta G < 0$, so the second law is
satisfied. Coupling is bookkeeping, not magic.

\textbf{The two requirements.} The reactions must \textbf{share a
common species} --- here \ce{P_i}, which cancels when you add them ---
and the favourable $\Delta G^{\circ}$ must be \textbf{larger in
magnitude} than the unfavourable one.

\textbf{The industrial version} appeared in Ladder 10: iron oxide
reduction on its own is hugely unfavourable
($\Delta G^{\circ} = \SI{+742}{\kilo\joule}$), but coupled to carbon
monoxide oxidation it comes out at \SI{-29.4}{\kilo\joule}. A blast
furnace is a coupled reaction the size of a building.
\end{workedexample}

\begin{yourturn}[YOUR TURN 15 --- four questions]
\begin{enumerate}[leftmargin=*,itemsep=0.5em]
  \item Reaction A has $\Delta G^{\circ} = \SI{+25}{\kilo\joule}$ and
        reaction B has $\SI{-40}{\kilo\joule}$. Find $\Delta G^{\circ}$
        for the coupled process. \workspace[1.8cm]
  \item Is the coupled process favoured? \workspace[1.4cm]
  \item What two conditions must hold for coupling to work?
        \workspace[1.8cm]
  \item Does coupling violate the second law? Explain.
        \workspace[1.8cm]
\end{enumerate}
\selfcheck{(a) \SI{-15}{\kilo\joule} \quad (b) yes \quad (c) a shared
species, and the favourable $\Delta G$ must be larger \quad (d) no ---
the total $\Delta G$ is still negative}
\keyonly{\begin{apcanswer}(a) $+25 - 40 =
\mathbf{\SI{-15}{\kilo\joule}}$. (b) \textbf{Yes} ---
$\Delta G^{\circ}$ is negative. (c) The two reactions must
\textbf{share a common species} so that adding them produces a sensible
overall equation, and the favourable reaction's $\Delta G^{\circ}$ must
be \textbf{larger in magnitude} than the unfavourable one's. (d)
\textbf{No.} The uphill reaction is paid for by a larger downhill one,
and the combined process still has $\Delta G < 0$ --- so
$\Delta S_{\text{univ}}$ still increases. Coupling is bookkeeping, not a
loophole.\end{apcanswer}}
\end{yourturn}

% =====================================================================
\section*{\sffamily Mixed set \textbullet\ ten questions, no labels}

\begin{apcnote}[Why this section exists]
Every question so far arrived with its ladder attached. These ten do
not --- and in this chapter the first decision is usually \emph{which
equation}: $\Delta H - T\Delta S$, formation values, $-RT\ln K$, or
$\Delta G^{\circ} + RT\ln Q$.
\end{apcnote}

\begin{yourturn}[MIXED SET --- first five]
\begin{enumerate}[leftmargin=*,itemsep=0.5em,label=(\alph*)]
  \item Give the sign of $\Delta S$ for
        $\ce{2H2O2(l) -> 2H2O(l) + O2(g)}$. \workspace[1.6cm]
  \item $\Delta H = \SI{-80.0}{\kilo\joule}$, $\Delta S =
        \SI{-200.0}{\joule\per\kelvin}$. Find $\Delta G$ at
        \SI{298}{\kelvin}. \workspace[2.0cm]
  \item For that reaction, find the crossover temperature and say
        which side is favoured. \workspace[2.0cm]
  \item $K = \num{2.0e-5}$ at \SI{298}{\kelvin}
        ($RT = \SI{2478}{\joule\per\mole}$). Find $\Delta G^{\circ}$.
        \workspace[2.0cm]
  \item A reaction has $\Delta G^{\circ} = \SI{-150}{\kilo\joule}$ yet
        nothing observable happens. Explain. \workspace[1.8cm]
\end{enumerate}
\selfcheck{(a) positive \quad (b) \SI{-20.4}{\kilo\joule} \quad (c)
\SI{400}{\kelvin}, favoured below \quad (d)
\SI{+26.8}{\kilo\joule} \quad (e) kinetic control}
\keyonly{\begin{apcanswer}\textbf{(a)} \textbf{Positive} --- a gas is
produced where there was none. \textbf{(b)} $\Delta G = -80.0 -
(298)(-0.2000) = -80.0 + 59.6 =
\mathbf{\SI{-20.4}{\kilo\joule}}$ --- note the double negative.
\textbf{(c)} $T = \num{-80000}/(-200.0) = \mathbf{\SI{400}{\kelvin}}$;
$\Delta S$ is negative, so it is favoured \textbf{below}
\SI{400}{\kelvin}. \textbf{(d)} $\Delta G^{\circ} = -(2478)
\ln(\num{2.0e-5}) = -(2478)(-10.82) = \SI{26800}{\joule} =
\mathbf{\SI{+26.8}{\kilo\joule}}$. \textbf{(e)} It is
\textbf{thermodynamically favoured but kinetically controlled} --- the
activation energy is large, so the rate is negligible.
$\Delta G$ describes extent, not speed.\end{apcanswer}}
\end{yourturn}

\begin{yourturn}[MIXED SET --- last five]
\begin{enumerate}[leftmargin=*,itemsep=0.5em,label=(\alph*)]
  \item Using $S^{\circ}$ (\si{\joule\per\mole\per\kelvin}):
        \ce{N2O4(g)} $= 304.3$, \ce{NO2(g)} $= 240.1$, find
        $\Delta S^{\circ}$ for $\ce{N2O4(g) -> 2NO2(g)}$.
        \workspace[2.0cm]
  \item A reaction has $\Delta H > 0$ and $\Delta S < 0$. At what
        temperatures is it favoured? \workspace[1.6cm]
  \item Reaction A ($\Delta G^{\circ} = \SI{+18}{\kilo\joule}$) is
        coupled to B ($\SI{-45}{\kilo\joule}$). Find the coupled
        $\Delta G^{\circ}$. \workspace[1.8cm]
  \item $\Delta G^{\circ} = 0$ for a reaction. What is $K$?
        \workspace[1.4cm]
  \item A salt dissolves endothermically yet the process is favoured.
        Which term is responsible, and why? \workspace[1.8cm]
\end{enumerate}
\selfcheck{(a) \SI{+175.9}{\joule\per\kelvin} \quad (b) never \quad
(c) \SI{-27}{\kilo\joule} \quad (d) $K = 1$ \quad (e) $T\Delta S$ ---
lattice to free ions}
\keyonly{\begin{apcanswer}\textbf{(a)} $2(240.1) - 304.3 = 480.2 -
304.3 = \mathbf{\SI{+175.9}{\joule\per\kelvin}}$ --- positive, as one
mole of gas becoming two requires. \textbf{(b)} \textbf{Never.} Both
terms of $\Delta H - T\Delta S$ are positive at every temperature.
\textbf{(c)} $+18 - 45 = \mathbf{\SI{-27}{\kilo\joule}}$ --- favoured.
\textbf{(d)} $\mathbf{K = 1}$, since $\ln K = 0$. \textbf{(e)} The
\textbf{entropy term $T\Delta S$}. Dissolving takes ions from fixed
lattice positions to freely moving hydrated ions, a large entropy
increase, and $T\Delta S$ outweighs the positive $\Delta H$
(Ladder 14).\end{apcanswer}}
\end{yourturn}

% =====================================================================
\section*{\sffamily Mastery tracker}

Tick a row only if \textbf{all four} YOUR TURN questions were right on
the first attempt. Every row is assessed --- CED topics 9.1 to 9.7
carry no exclusion statements at all, so nothing here is optional.

\renewcommand{\arraystretch}{1.25}
\noindent\begin{tabular}{@{}c p{5.7cm} c p{4.9cm}@{}}
\toprule
\textbf{First try?} & \textbf{Skill} & \textbf{Ladder} &
  \textbf{If not, re-read\ldots}\\
\midrule
$\square$ & what entropy is & 1 & arrangements, not mess \\
$\square$ & predicting the sign of $\Delta S$ & 2 & count gas moles first \\
$\square$ & the third law, absolute $S$ & 3 & elements are not zero \\
$\square$ & calculating $\Delta S^{\circ}$ & 4 & predict the sign first \\
$\square$ & the second law & 5 & only the sum must rise \\
$\square$ & $\Delta S_{\text{surr}}$ from $\Delta H$ & 6 & $-\Delta H/T$ \\
$\square$ & Gibbs free energy & 7 & convert kJ and J \\
$\square$ & the four sign cases & 8 & $T$ scales entropy only \\
$\square$ & the crossover temperature & 9 & $T = \Delta H/\Delta S$ \\
$\square$ & $\Delta G^{\circ}$ from formation & 10 & elements ARE zero here \\
$\square$ & thermodynamic vs kinetic & 11 & extent is not speed \\
$\square$ & $\Delta G^{\circ}$ and $K$ & 12 & $-RT\ln K$, joules \\
$\square$ & $\Delta G$ versus $\Delta G^{\circ}$ & 13 & $Q$ versus $K$ \\
$\square$ & free energy of dissolution & 14 & entropy usually wins \\
$\square$ & coupled reactions & 15 & the values add \\
\midrule
$\square$ & \textbf{mixed set} (8 of 10) & --- & whichever it exposed \\
\bottomrule
\end{tabular}

\begin{apcnote}[Scoring yourself honestly]
15/15 and the mixed set: you have Unit 9's thermodynamics half, and
Chapter 18 will lean on Ladders 12 and 13 constantly --- cell potential
is free energy wearing different units.

11--14: check where the misses fall. Ladders 1--6 are conceptual and
repair by re-reading. Ladders 7--10 are the calculation core, and the
overwhelming majority of errors there are the \textbf{kJ/J unit
mismatch} rather than any misunderstanding. If that is your pattern,
the fix is mechanical: write the units next to every number before
substituting.

10 or fewer: the likeliest single cause is treating $\Delta G$ as a
statement about \emph{speed}. Re-read Ladders 7 and 11 together, then
Ladder 12 --- ``favoured'' means the equilibrium lies toward products,
and nothing more. Everything else in this chapter is arithmetic on top
of that one distinction.

\textbf{One thing to carry into Chapter 18.} $\Delta G^{\circ} =
-nFE^{\circ}$ is coming, and it is Ladder 12 with a different constant.
If $\Delta G^{\circ} = -RT\ln K$ is solid, electrochemistry's
thermodynamics is already half learned.
\end{apcnote}

\end{document}
